Si consideri un cluster distribuito i cui nodi ospitano localmente parte dei dataset del sistema. Dall'esterno, degli utenti, detti emittenti, sottomettono richieste di esecuzione di task dichiarando i dataset e i requisiti necessari all'esecuzione, senza alcuna visibilità né sulla struttura del cluster né su dove i dati risiedano. Scegliere il nodo su cui far girare un task, però, non può ridursi a una questione di risorse disponibili. Sia i dati coinvolti sia il task stesso possono portare vincoli propri che vanno rispettati, e il loro rispetto richiede un insieme di politiche definite a livello di sistema, pensate per governare proprio la distribuzione delle computazioni che li riguardano. Un dataset, ad esempio, può essere vincolato a una certa area geografica al di fuori della quale non è trasferibile, richiedere garanzie minime al nodo che lo ospita, oppure appartenere a un contesto organizzativo in conflitto di interessi con altri; un task può imporre requisiti dello stesso tipo anche a prescindere dai dataset che utilizza.

Da qui nasce l'obiettivo della tesi, ovvero definire un modello formale di scheduling \emph{policy-driven} capace di governare questa decisione, per poi costruirne il prototipo. L'approccio riprende quello dei modelli formali proposti in letteratura per l'allocazione sicura dei dati in scenari multicloud, con la differenza che qui a essere distribuite non sono le informazioni ma le computazioni che le richiedono. Poiché la collocazione fisica dei dati si assume corretta fin dall'inizio, il modello non si occupa di dove i dati stiano fisicamente, ma valuta se, e a quale costo, un dataset possa essere spostato o acceduto da remoto verso il nodo scelto per l'esecuzione. L'assegnazione di un task si articola quindi in due fasi: una prima di filtraggio, in cui il task viene controllato nella sua ammissibilità complessiva e poi confrontato uno per uno con i nodi candidati, scartando quelli che violano i suoi vincoli o quelli dei dataset richiesti, e una seconda in cui, tra i nodi sopravvissuti al filtro, il migliore viene scelto tramite una funzione di ottimizzazione.

Le due fasi seguono innanzitutto un criterio di \emph{sovranità dei dati}, tema ormai centrale nella gestione di dati soggetti a giurisdizioni diverse, per cui le località dei nodi sono organizzate in aree geografiche di competenza e un task può essere eseguito solo entro l'area compatibile sia con i propri vincoli sia con quelli di ogni dataset che utilizza. Si tiene conto poi della \emph{località dei dati}, principio altrettanto consolidato nel calcolo distribuito, spingendo la computazione verso i nodi che già possiedono i dataset richiesti, così da ridurre quanto va trasferito o acceduto da remoto; del resto, per i dataset più sensibili e vulnerabili, anche un trasferimento controllato può risultare inaccettabile, ed è quindi ammessa solo l'esecuzione sul nodo che già li ospita. Inoltre, per valutare in modo uniforme le garanzie offerte da un nodo, ciascuno viene classificato rispetto a un insieme di proprietà astratte organizzate su livelli, ricavate a partire dai suoi attributi concreti, tecnici e non. A ogni nodo viene così assegnato, proprietà per proprietà, il livello più alto che i suoi attributi permettono di raggiungere, e sia i task sia i dataset esprimono i propri requisiti negli stessi termini, indicando semplicemente il livello minimo che pretendono dal nodo di esecuzione.

C'è poi un ultimo aspetto, forse il più delicato, ossia i conflitti di interesse. Eseguire un task può lasciare sul nodo tracce residue dei dati elaborati, e in linea di principio un task successivo, magari sottomesso da un concorrente diretto, potrebbe arrivare a quelle tracce. Per evitarlo, si adotta una politica di \emph{Chinese Wall}, riprendendo il classico modello di controllo degli accessi per i conflitti di interesse: ogni nodo accumula una memoria dei contesti organizzativi dei dati trattati dai task eseguiti in precedenza, e un nuovo task non può finire su un nodo la cui memoria è in conflitto con i contesti a cui il suo emittente è autorizzato ad accedere. A differenza del modello originale, pensato per soggetti umani la cui conoscenza non si può cancellare, qui il soggetto è un nodo che può essere materialmente ripulito; pertanto, una procedura periodica di sanificazione rimuove queste tracce e azzera la memoria, rendendo il nodo di nuovo disponibile per chiunque.

Il prototipo realizza tutto questo su un cluster Kubernetes, appoggiandosi ai meccanismi di estensione che la piattaforma mette già a disposizione, dal controllo di ammissione delle richieste fino allo scheduler vero e proprio, invece di costruire un orchestratore dedicato da zero.

La valutazione sperimentale si è concentrata sull'aspetto più significativo del modello, ovvero il comportamento dinamico introdotto dalla politica di Chinese Wall, che dipende dallo stato corrente del cluster. Ne è stato analizzato l'impatto sulla latenza di scheduling al variare del grado di conflitto tra i contesti e della frequenza di sanificazione dei nodi, evidenziando un compromesso tra disponibilità dei nodi e reattività dello scheduling.

In conclusione, il lavoro ha definito un modello formale che porta i vincoli di protezione dei dati direttamente dentro la decisione di scheduling, realizzandone un prototipo integrato in un orchestratore reale. Tra gli sviluppi futuri rientrano l'implementazione concreta del trasferimento fisico dei dati verso il nodo scelto, ad oggi valutato solamente in termini di costo e ammissibilità, e il raffinamento della procedura di sanificazione tramite l'adozione di più strategie combinate, per gestire in modo più efficace il compromesso appena descritto.